\documentclass[12pt]{article}
\usepackage[utf8]{inputenc}
\usepackage[T1]{fontenc}
\usepackage{hyperref}
\title{Kuhane Language Resources for K-12 Curriculum Development}
\author{Compiled by Claude Sonnet 5.5 (claude-sonnet-5-5) for Andr\'as Kornai}
\date{}
\begin{document}
\maketitle

\section{Introduction}
Kuhane (ISO 639-3: sbs; also Chikuhane, Chisubia, Subiya, Sesubiya) is a Bantu language spoken along the Zambezi in Namibia (East Caprivi), Botswana (Chobe) and Zambia (Western Province). Published speaker estimates vary widely, from about 31,000 to 40,000. The Endangered Languages Project rates it as vulnerable. It is used in most domains except official ones such as government, mass media and education. Silozi is the formal language of education and media in the area. Nothing found shows Kuhane being taught in schools. Little published material exists, and this dossier is short for that reason.

\section{Primary Educational Resources}
No primer, textbook or school reader was found. The only school-related items found are the orthography materials below.
\begin{itemize}
\item Ndana, Mabuta and Chebanne (2017), a Chiikuhane (Subiya) manual with orthography published through CASAS, Cape Town. Only the bibliographic record was seen. \url{https://glottolog.org/resource/reference/id/569862}
\item Baumbach's typescript ``tentative orthography of Subiya'' (Caprivi, undated, 198x), issued by the Namibian Ministry of Education and Culture. It is mentioned in the Endangered Languages Project record. \url{https://www.endangeredlanguages.com/lang/3544?hl=pt-br}
\end{itemize}

\section{Linguistic and Reference Materials}
\begin{itemize}
\item Wikipedia, Kuhane language: overview, location and classification. \url{https://en.wikipedia.org/wiki/Kuhane_language}
\item Endangered Languages Project record for Kuhane, with vitality, speaker numbers and bibliography. \url{https://www.endangeredlanguages.com/lang/3544?hl=pt-br}
\item Omniglot, Kuhane (Chikuhane / Chisubia): writing and sample text. \url{https://www.omniglot.com/writing/kuhane.htm}
\item Wikipedia, Subia people: ethnographic context. \url{https://en.wikipedia.com/wiki/Basubia}
\end{itemize}
No dictionary was found in the searches.

\section{Religious and Community Resources}
Omniglot reports a New Testament translation in Kuhane and gives the Lord's Prayer. The publisher and a link to the text were not confirmed. \url{https://www.omniglot.com/writing/kuhane.htm}

\section{Digital and Audio Resources}
No audio, video, corpus or software resources were found.

\section{Curriculum Development Context}
Silozi is the school and official language in the area. Literate Kuhane users in Botswana are reported to apply Setswana spelling conventions. A tentative orthography exists for Namibia, and the 2017 manual would be the first source to consult. A curriculum would need new primers, a dictionary and standardised orthography work, ideally with community and Namibian Ministry of Education involvement.
\end{document}